\subsection{从演化核到强子动量份额}
\label{sec:qcd-momentum-evolution-applications}

DIS 测量的是结构函数，而 DGLAP 演化作用在整组部分子分布上。要把矩阵演化同一个可直接解释的物理量联系起来，可以先问：提高分辨率以后，强子的纵向动量在夸克与胶子之间怎样重新分配？这一问题不需要先猜测 PDF 的形状。对已经建立的卷积方程取动量矩，就得到一个闭合的有限维动力学。

沿用 $f_{i/H}(x,\mu_{\rm F})$ 的算符定义，记
\begin{equation}
 \mathcal X_q(\mu_{\rm F})\equiv\int_0^1\dd x\,x\Sigma(x,\mu_{\rm F}),
 \qquad
 \mathcal X_g(\mu_{\rm F})\equiv\int_0^1\dd x\,x f_{g/H}(x,\mu_{\rm F}),
 \label{eq:qcd-momentum-fractions}
\end{equation}
其中 $\Sigma=\sum_{f=1}^{n_f}(f_{f/H}+f_{\bar f/H})$，$n_f$ 是这一尺度区间内的活跃夸克味数。两份额均无量纲；在只含 QCD 部分子的动量和规则中，$\mathcal X_q+\mathcal X_g=1$。若同时演化光子 PDF，总和必须加上光子动量，以下二分量闭合不再适用。

取重整化尺度等于因子化尺度，并在固定 $n_f$、无质量领先阶演化下定义无量纲常数
\begin{equation*}
 A_q\equiv\frac43C_F,\qquad A_g\equiv\frac23n_fT_F.
\end{equation*}
由\eqnrefs{eq:dglap-general-evolution-dp68} 得到
\begin{equation}
 \frac{\dd}{\dd\ln\mu_{\rm F}^2}
 \begin{pmatrix}\mathcal X_q\\\mathcal X_g\end{pmatrix}
 =\frac{\alpha_s}{2\pi}
 \begin{pmatrix}-A_q&A_g\\A_q&-A_g\end{pmatrix}
 \begin{pmatrix}\mathcal X_q\\\mathcal X_g\end{pmatrix}
 +\order(\alpha_s^2).
 \label{eq:qcd-momentum-matrix}
\end{equation}
每一列的和为零，表示从一种部分子流出的动量必定流入另一种。零本征值保证总动量守恒；另一个本征值 $-(A_q+A_g)$ 控制偏离固定比例的衰减。这并不表示微观辐射停止，而是表示两个方向的平均动量输运达到平衡。

以 $\mu_{\rm F0}$ 表示初始因子化动量尺度，避免与真空磁导率 $\mu_0$ 混淆。结合一圈强耦合运行，领先对数解为
\begin{equation}
 \begin{aligned}
 \mathcal X_q(\mu_{\rm F})
 &=\mathcal X_q^*+
 [\mathcal X_q(\mu_{\rm F0})-\mathcal X_q^*]
 \left[\frac{\alpha_s(\mu_{\rm F})}{\alpha_s(\mu_{\rm F0})}\right]^{d_{\rm mom}},\\
 \mathcal X_q^*&=\frac{A_g}{A_q+A_g},\qquad
 \mathcal X_g^*=\frac{A_q}{A_q+A_g},\qquad
 d_{\rm mom}=\frac{2(A_q+A_g)}{\beta_0}.
 \end{aligned}
 \label{eq:qcd-momentum-solution}
\end{equation}
这里 $\beta_0=(11C_A-4T_Fn_f)/3$。对于渐近自由理论，$\beta_0>0$；提高尺度使 $\alpha_s$ 减小，偏离固定比例的分量因而减小。真实 QCD 取 $C_F=4/3,T_F=1/2$ 后，固定比例成为
\begin{equation*}
 \mathcal X_q^*=\frac{3n_f}{16+3n_f},\qquad
 \mathcal X_g^*=\frac{16}{16+3n_f}.
\end{equation*}
这是固定味数的 QCD 渐近结果。穿过重夸克阈值必须分段演化并作匹配；任意有限尺度上的真实份额还取决于 PDF 初值，不能直接用固定比例替代实验输入。

\begin{exbox}[由给定初值计算动量输运]
在不穿过重味阈值的五味教学模型中，取
\begin{equation*}
 c\mu_{\rm F0}=\SI{10}{\giga\electronvolt},\quad
 c\mu_{\rm F}=\SI{100}{\giga\electronvolt},\quad
 \alpha_s(\mu_{\rm F0})=0.180,\quad
 \mathcal X_q(\mu_{\rm F0})=0.600.
\end{equation*}
一圈演化给出
\begin{align*}
 \alpha_s(\mu_{\rm F})
 &=\frac{0.180}{1+(23/3)(0.180)\ln(100)/(4\pi)}\simeq0.1195,\\
 \mathcal X_q^*&=\frac{15}{31}\simeq0.4839,\qquad
 d_{\rm mom}=\frac{62}{69},\\
 \mathcal X_q(\mu_{\rm F})
 &\simeq0.4839+(0.600-0.4839)(0.1195/0.180)^{62/69}
 \simeq0.5643,\\
 \mathcal X_g(\mu_{\rm F})&\simeq0.4357.
\end{align*}
约 $3.57\%$ 的强子总动量从夸克扇区转移到胶子扇区。这组数值来自给定初值的演化算例；要用于质子预测，还需由数据确定 PDF 初值。微扰截断、重味匹配和 PDF 初值的不确定性应分别报告。
\end{exbox}

\begin{figure}[htbp]
 \centering
 \includegraphics[width=0.73\textwidth]{qcd_momentum_evolution.pdf}
 \caption{五味领先阶 QCD 中的夸克与胶子动量份额。曲线由正文给定初值演化得到，虚线为固定味数平衡比例；二者之和始终为一。}
 \label{fig:qcd-momentum-evolution}
\end{figure}
\figref{fig:qcd-momentum-evolution} 中趋近固定比例的速度很慢，因为控制变量是运行耦合的幂，等价于硬尺度对数的幂。仅凭两条份额曲线接近，不能推断它们的完整 $x$ 分布也接近。

\begin{derivation}[动量矩的闭合与对角化]
在矩积分收敛的条件下，交换积分次序并令 $x=zu$，
\begin{align*}
 \int_0^1\dd x\,x[P_{ij}\otimes f_j](x)
 &=\int_0^1\dd z\,P_{ij}(z)\int_0^z\frac{\dd x}{z}\,x f_j(x/z)\\
 &=\left[\int_0^1\dd z\,zP_{ij}(z)\right]
   \left[\int_0^1\dd u\,u f_j(u)\right].
\end{align*}
用前文含虚修正的完整分裂核，四个矩为
\begin{align*}
 P_{qq,2}^{(0)}&=-\frac43C_F=-A_q,&
 P_{gq,2}^{(0)}&=\frac43C_F=A_q,\\
 2n_fP_{qg,2}^{(0)}&=\frac23n_fT_F=A_g,&
 P_{gg,2}^{(0)}&=-\frac23n_fT_F=-A_g.
\end{align*}
这里胶子馈入味单态的 $2n_f$ 来自每个味的夸克和反夸克；漏掉因子 $2$ 会同时破坏动量守恒和固定比例。代入即得\eqnrefs{eq:qcd-momentum-matrix}。

在总动量为一的子空间令 $\delta\mathcal X_q=\mathcal X_q-\mathcal X_q^*$，则
\begin{align*}
 \frac{\dd\delta\mathcal X_q}{\dd\ln\mu_{\rm F}^2}
 &=-\frac{\alpha_s}{2\pi}(A_q+A_g)\delta\mathcal X_q,\\
 \frac{\dd\ln|\delta\mathcal X_q|}{\dd\alpha_s}
 &=\frac{2(A_q+A_g)}{\beta_0\alpha_s}.
\end{align*}
对第二式积分便得\eqnrefs{eq:qcd-momentum-solution}；初始偏差为零时直接使用恒定解，无需求对数。若暂时冻结耦合，则偏差随 $\ln(\mu_{\rm F}^2/\mu_{\rm F0}^2)$ 指数衰减，这也检查了运行解指数的正负号。
\end{derivation}

\subsection{用束流能量扫描分离纵向结构函数}
\label{sec:dis-longitudinal-extraction}

动量份额属于 PDF 的积分性质，实验首先得到的仍是微分截面。由\eqnrefs{eq:dis-reduced-cross-section-dp31}，固定 $x_B,\mathcal Q^2$ 后，定义已知运动学权重 $w(y)=y^2/[1+(1-y)^2]$，约化截面就是关于 $w$ 的直线：截距为 $F_2$，斜率为 $-F_L$。两个束流能量分别对应 $y_1,y_2$，测得 $r_i=\sigma_r(y_i)$，则
\begin{equation}
 F_L=\frac{r_1-r_2}{w_2-w_1},\qquad
 F_2=\frac{w_2r_1-w_1r_2}{w_2-w_1},\qquad w_i=w(y_i).
 \label{eq:dis-two-energy-inversion}
\end{equation}
在忽略靶质量、纯光子交换的近似内，纵向与横向吸收截面非负要求 $0\leq F_L\leq F_2$。这提供了独立于直线拟合的物理检查。较大 $y$ 增强纵向灵敏度，但单纯提高统计量不能替代不同 $w_i$ 之间的间隔。

\begin{exbox}[纵向响应的数值提取与误差]
取同一 $(x_B,\mathcal Q^2)$ 上的两点 $y_1=0.2,y_2=0.8$，得到
$w_1=0.02439,w_2=0.61538$。若教学数据为 $r_1=0.49756,r_2=0.43846$，则
\begin{equation*}
 F_L\simeq\frac{0.05910}{0.59099}=0.1000,\qquad
 F_2=r_1+w_1F_L\simeq0.5000.
\end{equation*}
若两个截面具有标准差 $s_1,s_2$、协方差 $V_{12}$，线性误差传播给
\begin{equation*}
 \operatorname{Var}(F_L)
 =\frac{s_1^2+s_2^2-2V_{12}}{(w_2-w_1)^2}.
\end{equation*}
取独立误差 $s_1=s_2=0.005$，得到 $s(F_L)\simeq0.0120$。若两点的 $y$ 很接近，分母趋于零，反演便不稳定；共享束流归一化误差则应进入 $V_{12}$，不能按两份独立噪声处理。实际分离还需保持相同的 $(x_B,\mathcal Q^2)$ 分箱，并校正辐射、接收度以及非光子交换贡献。
\end{exbox}

\begin{derivation}[有限靶质量的纵向投影]
前文 $F_L=F_2-2x_BF_1$ 是领先靶质量约定。若保留靶质量 $M_H$，从\eqnrefs{eq:dis-transverse-tensor-basis-dp68} 得到
\begin{equation*}
 \widetilde P^2=M_H^2c^2+\frac{(P\cdot q)^2}{\mathcal Q^2},\qquad
 \varepsilon_L^\mu=\frac{\widetilde P^\mu}{\sqrt{\widetilde P^2}},\qquad
 \varepsilon_L^2=1,\quad q\cdot\varepsilon_L=0.
\end{equation*}
类空虚光子正交子空间中的这个类时向量不等同于实光子的横向极化。其强子张量投影为
\begin{align*}
 W_L&=\varepsilon_L^\mu W_{\mu\nu}\varepsilon_L^\nu
 =-F_1+\frac{\widetilde P^2}{P\cdot q}F_2,\\
 2x_BW_L&=\left(1+\frac{4x_B^2M_H^2c^2}{\mathcal Q^2}\right)F_2-2x_BF_1.
\end{align*}
所以保留有限靶质量时，物理纵向组合应记为
\begin{equation}
 F_L^{\rm phys}\equiv
 \left(1+\frac{4x_B^2M_H^2c^2}{\mathcal Q^2}\right)F_2-2x_BF_1.
 \label{eq:dis-longitudinal-target-mass}
\end{equation}
当 $4x_B^2M_H^2c^2/\mathcal Q^2\ll1$ 时才恢复前述直线分离的领先幂形式。这个运动学修正与结构函数内部的高扭曲动力学修正不同：只修改纵向定义，并不能把全部有限 $\mathcal Q$ 效应一并纳入。
\end{derivation}
